\chapter{GIARDINAGGIO CELESTE}\label{giardinaggio-celeste}

\hfill\break

Quello fu l'inverno delle torri di lancio.

Erano state erette nuove piattaforme non solo a Canaveral ma anche in
tutto il Sud-Ovest deserto, nel sud della Francia e nell'Africa
equatoriale, a Jiuquan e Xichang in Cina e a Bayqoñyr e Svobodnyj in
Russia: torri per i lanci di semi marziani e torri più grandi per le
cosiddette Grosse Cataste, gli enormi razzi vettori che avrebbero
portato volontari umani su un Marte parzialmente abitabile se la
rudimentale terraformazione fosse riuscita. Quell'inverno le torri
crebbero come foreste di ferro e acciaio: lussureggianti, rigogliose,
radicate nel cemento e annaffiate con cisterne di denaro federale.

Per certi versi i primi razzi con i semi erano meno spettacolari delle
strutture di lancio costruite per sostenerli. Erano razzi vettori da
catena di montaggio prodotti in serie da vecchi stampi di Titan e Delta,
non avevano neanche un microchip in più di quelli strettamente
necessari. Ce n'era una quantità sorprendente a popolare le piattaforme
mentre l'inverno avanzava verso la primavera, astronavi come capsule di
pioppo nero, pronte a portare la vita dormiente in un suolo lontano e
sterile.

In un certo senso, fu primavera anche nel sistema solare, o per lo meno
un'estate indiana prolungata. La zona abitabile del sistema solare si
spostava verso l'esterno, poiché il Sole stava esaurendo il suo nucleo
di elio, e cominciava a includere Marte così come alla fine avrebbe
incluso Ganimede, la luna di Giove, un altro potenziale obiettivo per
una terraformazione successiva. Su Marte, dopo milioni di estati calde,
avevano iniziato a sublimare nell'atmosfera quantità enormi di
CO\textsubscript{2} congelata e ghiaccio d'acqua. All'inizio dello Spin,
la pressione atmosferica di Marte a livello del suolo era di circa otto
millibar, rarefatta come l'aria che si trova a cinque chilometri sopra
la cima del monte Everest. In seguito, anche senza l'intervento umano,
il pianeta aveva raggiunto un clima equivalente a quello di una vetta
artica immersa in diossido di carbonio gassoso, cioè mite per gli
standard marziani.

Ma la nostra intenzione era spingerci oltre. Volevamo aggiungere
ossigeno all'aria del pianeta, inverdire le sue pianure, creare stagni
lì dove il ghiaccio sotto la superficie si scioglieva periodicamente ed
eruttava sotto forma di geyser di vapore o depositi semiliquidi di fango
tossico.

Durante l'inverno delle torri di lancio fummo pericolosamente ottimisti.

ooo

Il tre marzo, poco prima della prima ondata programmata di lanci di
semi, Carol Lawton mi chiamò a casa e mi disse che mia madre aveva avuto
un grave ictus e non c'erano speranze che sopravvivesse.

Presi accordi perché un medico del posto mi sostituisse alla Perielio,
poi andai in macchina a Orlando e prenotai il primo volo diurno per
Washington D.C.

Carol venne a prendermi al Reagan International. Sembrava sobria.
Spalancò le braccia e io la strinsi, quella donna che durante gli anni
in cui avevo vissuto nella sua proprietà mi aveva dimostrato solo
un'indifferenza confusa. Poi si ritrasse e mi poggiò le mani tremanti
sulle spalle. «Mi dispiace tanto, Tyler.»

«È ancora viva?»

«Tiene duro. C'è una macchina che mi aspetta. Possiamo parlare lungo il
tragitto.»

La seguii fino a un'auto che doveva essere stata mandata dallo stesso
E.D., una limousine nera con adesivi federali. L'autista parlò a
malapena mentre infilava nel bagagliaio le mie valigie, quando lo
ringraziai sollevò appena il cappello e salì sul sedile di guida
meticolosamente isolato dal lussuoso abitacolo per i passeggeri. Si
diresse verso l'ospedale della George Washington University senza
bisogno di chiederglielo.

Carol era più magra di come la ricordassi, un uccellino sulla
tappezzeria in pelle. Prese un fazzoletto di cotone dalla sua borsetta e
si tamponò gli occhi. «Che assurdità piangere in questo modo. Ieri ho
perso le lenti a contatto. Potrei dire che le ho piante via, te
l'immagini? Ci sono cose che si danno per scontate. Una di queste per me
era avere tua madre in casa che teneva le cose in ordine o semplicemente
saperla vicina, lì in fondo al prato. Mi svegliavo di notte - non dormo
profondamente, e forse non ne sei sorpreso - mi svegliavo con la
sensazione che il mondo fosse fragile e io potessi sprofondarvi dentro,
sprofondare nel pavimento e continuare a cadere per sempre. Poi pensavo
a lei laggiù nella Piccola Casa, che dormiva profondamente. Dormire
profondamente. Era come una prova giudiziaria. Reperto A, Belinda
Dupree, la dimostrazione che la tranquillità d'animo è possibile. Era la
colonna della famiglia, Tyler, che tu lo sapessi o no.»

Supponevo di averlo saputo. Eravamo stati davvero una sola famiglia,
anche se da bambino avevo notato soprattutto la differenza tra le due
proprietà: la mia casa, modesta ma tranquilla, e la Grande Casa, dove i
giocattoli erano più costosi e le discussioni più violente.

Chiesi se E.D. fosse andato in ospedale.

«E.D.? No. È impegnato. Pare che mandare astronavi su Marte richieda
tantissime cene in centro. So che anche Jason resta in Florida per
questo, ma credo che lui si occupi del lato pratico della questione, se
\emph{esiste} un lato pratico, mentre E.D. è più un mago da
palcoscenico, tira fuori i soldi da vari cappelli. Ma sono certa che
vedrai E.D. al funerale.» Trasalii, e lei mi guardò con aria di scusa.
«Se e quando dovesse accadere. Ma i dottori dicono\ldots»

«Che non ci sono speranze che si rimetta.»

«Sta morendo. Sì. Da medico a medico. Te lo ricordi, Tyler? Una volta
esercitavo anch'io. Ai tempi in cui ero capace di una cosa del genere. E
ora anche tu sei un medico che esercita per conto suo. Mio Dio.»

Apprezzai la sua schiettezza. Forse comparve con l'improvvisa sobrietà.
Eccola tornata nel mondo illuminato a giorno che evitava da vent'anni,
orribile proprio come se lo ricordava.

Arrivammo all'ospedale della George Washington University. Carol si era
già presentata al personale infermieristico del reparto rianimazione,
quindi ci dirigemmo direttamente verso la stanza di mia madre. Quando
vidi Carol esitare sulla porta, le chiesi: «Entri?»

«Io\ldots{} no, non credo. Le ho già detto addio tante volte. Ho bisogno
di stare all'aria aperta, dove l'aria non puzza di disinfettante. Me ne
starò nel parcheggio a fumare una sigaretta con i barellieri. Ci
troviamo lì?»

Le dissi di sì.

Mia madre era incosciente, attaccata alle macchine, con un respiratore
che le regolava il respiro e sibilava mentre la sua gabbia toracica si
gonfiava e si rilassava. Aveva i capelli più bianchi rispetto a come me
la ricordavo io. Le carezzai la guancia, lei non reagì.

Per un qualche istinto strampalato di medico, le sollevai una palpebra,
con l'intento, credo, di controllare la dilatazione delle pupille. Ma
dopo l'ictus aveva avuto un'emorragia all'occhio. Era cremisi come un
pomodoro, arrossato dal sangue.

ooo

Me ne andai dall'ospedale con Carol ma rifiutai il suo invito a cena, le
dissi che mi sarei preparato qualcosa io stesso. Rispose: «Sono sicura
che in casa di tua madre ci sia qualcosa. Ma se vuoi stare nella Grande
Casa sei più che benvenuto. Anche se è un po' un casino in questi giorni
senza tua madre che tiene tutto sotto controllo. Sono certa che possiamo
improvvisare una camera degli ospiti passabile.»

La ringraziai spiegandole che preferivo stare dall'altra parte del
prato.

«Fammi sapere se cambi idea.» Fissò la Piccola Casa dal vialetto
ghiaioso come se dopo anni la vedesse con chiarezza per la prima volta.
«Hai ancora la chiave con te?»

«Ce l'ho.»

«Bene, allora. Ti lascio andare. L'ospedale ha i numeri di tutti e due
per avvisarci se le sue condizioni dovessero cambiare.» E mi
riabbracciò, poi salì le scale del portico con una certa risolutezza ma
senza slancio; si intuiva che aveva rimandato già abbastanza a lungo il
suo bisogno di bere.

Aprii la porta della casa di mia madre. ``Sua più che mia'' pensai,
sebbene la mia presenza non vi fosse stata cancellata. Quando ero
partito per l'università avevo svuotato la mia cameretta e impacchettato
tutto quello che ritenevo importante, ma mia madre aveva lasciato il
letto e riempito gli spazi vuoti (gli scaffali di pino, il davanzale)
con piante in vaso, che in sua assenza si erano rapidamente seccate; le
annaffiai. Anche il resto della casa era in ordine. Diane una volta
aveva definito «lineare» il modo in cui mia madre teneva la casa, e con
questo credo che intendesse ordinato ma non ossessivo. Controllai il
soggiorno, la cucina, diedi un'occhiata alla sua camera da letto. Non
tutto era al suo posto.

Ma tutto \emph{aveva} un posto.

All'imbrunire chiusi le tende e accesi tutte le luci di ogni stanza, una
quantità di luci che mia madre non aveva mai ritenuto opportuno
accendere, una dichiarazione contro la morte. Mi chiesi se Carol avesse
notato il bagliore oltre la spianata resa marrone dal clima invernale e,
nel caso, se lo trovasse confortante o allarmante.

Quella sera E.D. tornò a casa verso le nove e fu abbastanza gentile da
bussare alla porta e mostrarsi comprensivo. Sembrava a disagio sotto la
luce del portico, con l'abito di buon taglio in disordine. Nel gelo
della sera il suo respiro fumava. Si tastò le tasche, il petto e il
fianco, senza rendersene conto, come se avesse dimenticato qualcosa o
semplicemente non sapesse cosa fare con le mani. «Mi dispiace, Tyler.»

Le sue condoglianze mi sembrarono volgarmente premature, come se la
morte di mia madre fosse non solo inevitabile ma anche un fatto
incontrovertibile. L'aveva già cancellata. ``Ma lei respira ancora,''
pensai, ``o almeno fa entrare ancora ossigeno dentro di sé, a chilometri
di distanza, da sola nella sua stanza al George Washington.'' «Grazie
per le sue parole, signor Lawton.»

«Ti prego, Tyler, dammi del tu. Lo fanno tutti gli altri. Jason mi dice
che stai facendo un buon lavoro, laggiù alla Perielio, in Florida.»

«I miei pazienti sembrano soddisfatti.»

«Fantastico. Ogni contributo conta, anche se piccolo. Senti, ma è stata
Carol a farti stare qui? Perché abbiamo una stanza degli ospiti pronta,
se vuoi.»

«Sto benissimo dove sono.»

«Va bene. Lo capisco. Bussa senza problemi se hai bisogno di qualcosa,
va bene?»

Tornò indietro attraversando il prato marrone con passo tranquillo.
C'era stato un gran parlare, nella stampa e nella famiglia Lawton, della
genialità di Jason, ma rammentai a me stesso che anche E.D. avrebbe
potuto rivendicare quella qualità. Aveva messo a frutto la laurea in
ingegneria e il talento per gli affari in un'importante società e aveva
cominciato a vendere banda larga per le telecomunicazioni grazie agli
aerostati quando l'Americom e l'AT\&T erano rimaste immobili di fronte
allo Spin come cervi spaventati. Non gli mancava l'intelligenza di
Jason, ma senz'altro la sua arguzia e la profonda curiosità per
l'universo fisico. E forse un pizzico della sua umanità.

Poi fui di nuovo solo, in quella che era stata la mia casa ma che ormai
non lo era più. Sedetti sul divano e per un attimo mi stupii di quanto
poco fosse cambiata quella stanza. Prima o poi sarebbe toccato a me
disporre del contenuto della casa, un compito che a malapena riuscivo a
immaginare, un compito più difficile, più illogico del lavoro di
coltivare la vita su un altro pianeta. E forse fu perché stavo meditando
su quell'atto di decostruzione che notai uno spazio vuoto sull'ultimo
ripiano dell'étagère accanto alla TV.

Lo notai perché, per quanto ne sapevo, in tutti gli anni vissuti lì il
ripiano in alto era stato spolverato solo superficialmente. L'ultimo
ripiano era l'attico della vita di mia madre. Avrei potuto recitare a
memoria e a occhi chiusi l'ordine delle cose poggiate sopra, solo
raffigurandole nella mente: i suoi annuari delle superiori (la scuola
superiore Martell a Bingham, nel Maine, 1975, 76, 77, 78); la guida
dello studente della Berkeley, 1982; un Buddha di giada per reggilibro;
il suo diploma in una cornice di plastica con il sostegno; un
raccoglitore marrone a soffietto in cui conservava certificato di
nascita, passaporto e documenti fiscali; e, sorrette da un altro Buddha
verde, tre malconce scatole di scarpe New Balance con le etichette
{RICORDI} ({SCUOLA}), {RICORDI} ({MARCUS}) e {CIANFRUSAGLIE}.

Ma quella sera il secondo Buddha di giada era spostato e mancava la
scatola contrassegnata con {RICORDI} ({SCUOLA}). Ipotizzai che l'avesse
tirata giù lei stessa, anche se non l'avevo vista da nessun'altra parte
in casa. Delle tre scatole, l'unica che aveva aperto abitualmente in mia
presenza era quella {CIANFRUSAGLIE}. Era stata riempita di locandine di
concerti e matrici di biglietti, ritagli lisi di giornale (fra cui i
necrologi dei suoi genitori), una spilla con la forma della goletta
\emph{Bluenose}, ricordo della luna di miele in Nuova Scozia, bustine di
fiammiferi di ristoranti e alberghi in cui era stata, bigiotteria, un
certificato di battesimo, persino una ciocca di capelli di quando ero
piccolo conservata in una striscia di carta cerata chiusa con una
spilla.

Tirai giù l'altra scatola, quella contrassegnata con {RICORDI}
({MARCUS}). Non ero mai stato particolarmente curioso di sapere di mio
padre, e mia madre ne parlava raramente, se non per sommi capi: un
bell'uomo, un ingegnere, un collezionista di jazz, il migliore amico di
E.D. al college, ma anche un gran bevitore che una notte, tornando a
casa dopo essere stato da un fornitore di dispositivi elettronici a
Milpitas, era rimasto vittima della sua stessa passione per le auto
veloci. Dentro la scatola da scarpe c'era una pila di lettere in buste
di pergamena su cui l'indirizzo era stato scritto con una grafia
asciutta e pulita che dedussi essere la sua. Le aveva spedite a Belinda
Sutton, il cognome da nubile di mia madre, a un indirizzo di Berkeley
che non riconobbi.

Sfilai una busta dal mucchio e la aprii, tirai fuori la carta ingiallita
e la spiegai.

Il foglio non aveva righe, ma la scrittura attraversava la pagina con
brevi parallele ordinate. ``\emph{Cara Bel},'' iniziava, e poi
proseguiva, ``\emph{ieri sera pensavo di averti detto tutto al telefono
	e invece non riesco a smettere di pensarti. Scrivendoti queste parole mi
	sembra di averti più vicina, anche se non così vicina come vorrei. Non
	come lo eravamo lo scorso agosto! Con la mente mi immergo nella memoria
	di quei giorni ogni notte in cui non posso sdraiarmi accanto a te}.''

C'era altro ma non lo lessi. Piegai la lettera e la infilai nella busta
ingiallita, chiusi la scatola e la rimisi al suo posto.

ooo

Al mattino bussarono alla porta. Risposi aspettandomi che fosse Carol o
qualche segretario che veniva dalla Grande Casa.

Ma non era Carol. Era Diane. Diane che indossava una gonna blu notte
lunga fino ai piedi e una camicetta a collo alto. Aveva le mani
intrecciate sotto il seno. Alzò lo sguardo verso di me, gli occhi
luccicavano. «Mi dispiace tanto. Sono venuta appena ho saputo.»

Ma troppo tardi. L'ospedale mi aveva chiamato dieci minuti prima.
Belinda Dupree era morta senza riprendere conoscenza.

ooo

Al servizio funebre E.D. era a disagio, parlò brevemente e non disse
nulla di significativo. Parlai io, parlò Diane; Carol voleva ma alla
fine non lo fece perché piangeva troppo o era troppo ubriaca per salire
sul pulpito.

L'elogio di Diane fu il più commovente, sentito e con le giuste pause.
Un catalogo delle gentilezze che mia madre aveva esportato attraverso il
prato come regali di una nazione più ricca e gentile. Ne fui grato. Al
confronto tutto il resto della cerimonia parve meccanico: facce quasi
familiari comparivano all'improvviso tra la folla per pronunciare omelie
e mezze verità, io li ringraziavo e sorridevo, li ringraziavo e
sorridevo, finché arrivò il momento del corteo funebre per raggiungere
la sepoltura.

ooo

Quella sera ci fu una funzione alla Grande Casa, un ricevimento dopo il
funerale in cui ricevetti le condoglianze dai soci in affari di E.D., di
cui non conoscevo nessuno sebbene alcuni avessero conosciuto mio padre,
e dai domestici della Grande Casa, il cui dolore era più autentico e più
difficile da sopportare.

Gli addetti al catering scivolavano tra la folla con bicchieri colmi di
vino su piatti d'argento; bevvi più di quanto avrei dovuto finché Diane,
che si era mescolata anche lei in silenzio tra gli ospiti, mi trascinò
via da un altro giro di ``mi spiace così tanto per la tua perdita'' e
disse: «Hai bisogno di aria.»

«Fa freddo.»

«Se continui a bere diventerai scontroso. Già ci sei quasi. Dai, Ty.
Solo per qualche minuto.»

Uscimmo sul prato. Il prato marrone di pieno inverno. Lo stesso in cui
quasi vent'anni prima avevamo assistito ai momenti iniziali dello Spin.
Camminammo lungo la circonferenza della Grande Casa; gironzolammo, in
realtà, malgrado la brezza rigida di marzo e la neve granulosa che
occupava ancora ogni spazio riparato o all'ombra.

Ci eravamo già detti tutte le cose ovvie. Avevamo parlato delle
reciproche esperienze: la mia professione, il trasferimento in Florida,
il mio lavoro alla Perielio; gli anni passati con Simon, in cui si erano
allontanati dall'NR per abbracciare un'ortodossia più moderata,
accettando l'Assunzione con fede e sacrificio di sé. («Non mangiamo
carne» mi aveva confidato. «Non indossiamo fibre artificiali.»)
Camminando accanto a lei, con la testa leggera, mi chiedevo se ai suoi
occhi fossi diventato volgare o ripugnante, se si fosse accorta del mio
alito che sapeva di stuzzichini al prosciutto e formaggio o della giacca
di cotone e poliestere che indossavo. Diane non era cambiata molto,
anche se era più magra di prima, forse più magra di quanto avrebbe
dovuto essere, la linea della mascella un po' dura sul colletto alto e
stretto.

Ero abbastanza sobrio da riuscire a ringraziarla per aver cercato di
farmi passare la sbornia.

«Anch'io avevo bisogno di allontanarmi» confessò. «E poi tutte quelle
persone che E.D. ha invitato\ldots{} Nessuna conosceva bene tua madre.
Neanche una. Stanno lì a parlare di disegni di legge di stanziamento o
stazza dei carichi. A concludere affari.»

«Forse E.D. pensa sia questo il modo migliore in cui renderle omaggio.
Dà sapore alla veglia funebre con le celebrità politiche.»

«È un modo generoso di interpretare la cosa.»

«Ti fa ancora arrabbiare.» ``E con che facilità'' pensai.

«E.D.? Certo. Anche se sarebbe più caritatevole perdonarlo\ldots{} e mi
pare che tu l'abbia già fatto.»

«Ho meno cose da perdonargli, non è mio padre.»

E dicendo questo non volevo intendere nulla. Ma ero ancora troppo
sensibile a quello che Jason mi aveva detto qualche settimana prima.
Quelle parole mi uscirono strozzate dalla gola, le riesaminai ancora
prima che mi venissero fuori e arrossii dopo averle dette. Diane mi
guardò a lungo senza capire; poi spalancò gli occhi in un'espressione di
rabbia mista a imbarazzo, così evidente che riuscii a riconoscerla anche
al fioco bagliore della luce del portico.

«Hai parlato con Jason» sibilò.

«Mi dispiace\ldots»

«Come funziona esattamente questa cosa? Ve ne state lì a oziare
prendendomi in giro?»

«Certo che no. Lui\ldots{} qualunque cosa abbia detto Jason è dovuta ai
farmaci.»

Un altro grottesco passo falso, e lei lo colse al volo: «Che farmaci?»

«Sono il suo medico generico. A volte gli prescrivo delle ricette. Che
differenza fa?»

«Tyler, che genere di farmaci può farti infrangere una promessa? Aveva
promesso che non ti avrebbe mai detto che\ldots» trasse un'altra
conclusione, «Jason è \emph{malato}? È per questo che non è venuto al
funerale?»

«È impegnato. Mancano pochi giorni ai primi lanci.»

«Ma lo stai curando per un motivo.»

«Per etica professionale non posso parlare dell'anamnesi di Jason» le
dissi, sapendo che questo avrebbe alimentato ancora di più i suoi
sospetti e che, nel tentativo di mantenere il segreto di Jason, lo avevo
in realtà svelato.

«Sarebbe proprio da lui, ammalarsi e non farlo sapere a nessuno di noi.
È proprio un tipo \emph{sigillato ermeticamente}\ldots»

«Forse dovresti prendere tu l'iniziativa. Chiamalo, qualche volta.»

«Pensi che non lo faccia? Anche questo te l'ha detto lui? Lo chiamavo
ogni settimana. Ma lui faceva tante moine inutili e non diceva mai
niente di importante. ``Come stai, sto bene, c'è qualche novità,
niente.'' Lui non vuole sentirmi, Tyler. Sta dalla parte di E.D.,
completamente. Sono motivo di imbarazzo per lui.» Fece una pausa. «A
meno che non sia cambiato qualcosa.»

«Non so cosa sia cambiato. Ma forse dovresti andare a trovarlo,
parlargli di persona.»

«E come potrei?»

Scrollai le spalle. «Prenditi un'altra settimana libera. Vieni in aereo
con me.»

«Hai detto che è impegnato.»

«Una volta iniziati i lanci, c'è solo da sedersi e aspettare. Puoi
venire a Canaveral con noi. Vedrai fare la Storia.»

«I lanci sono inutili» disse, ma sembrava che qualcuno le avesse
insegnato a dire così. Aggiunse: «Mi piacerebbe, ma non posso
permettermelo. Io e Simon ce la caviamo ma non siamo ricchi. Non siamo i
Lawton.»

«Te li presto io i soldi per il biglietto dell'aereo.»

«Sei un ubriaco generoso.»

«Dico sul serio.»

«Grazie, ma no, non potrei.»

«Pensaci.»

«Chiedimelo quando sarai sobrio.» E mentre salivamo i gradini del
portico, con la luce gialla che le incappucciava gli occhi, aggiunse:
«Qualunque cosa possa aver creduto un tempo\ldots{} qualunque cosa possa
aver detto a Jason\ldots»

«Non c'è bisogno che tu lo dica, Diane.»

«Lo so che E.D. non è tuo padre.»

La cosa interessante della sua ritrattazione fu il modo in cui la
pronunciò. In modo fermo, deciso. Come se in quel momento ne fosse
ancora più convinta. Come se avesse scoperto una verità diversa, una
chiave alternativa ai misteri dei Lawton.

ooo

Diane tornò alla Grande Casa e io decisi che non ce l'avrei fatta ad
affrontare altra gente che mi augurava buone cose. Entrai nella casa di
mia madre che sembrava senza aria e surriscaldata.

Il giorno dopo Carol mi disse che per sgombrare gli oggetti di mia
madre, cosa che lei definì ``organizzare'', potevo prendermi tutto il
tempo. «La Piccola Casa non andrà da nessuna parte» mi rassicurò.
«Prenditi un mese. Prenditi un anno.»

Potevo ``organizzare'' ogniqualvolta ne avessi avuto il tempo e quando
sarei stato più tranquillo.

In quel momento la tranquillità non era nemmeno all'orizzonte, ma la
ringraziai per la pazienza e passai la giornata a fare i bagagli per
tornare in volo a Orlando. Ero infastidito dal pensiero di dovermi
portare via qualcosa di mia madre e che lei avrebbe voluto che tenessi
un ricordo da conservare in una mia personale scatola di scarpe. Ma
cosa? Una delle sue statuine Hummel, che a lei piacevano tanto ma che a
me avevano sempre dato l'impressione di roba pacchiana e costosa? La
farfalla ricamata a punto croce sulla parete del soggiorno, la stampa
delle \emph{Ninfee} in una cornice fai da te?

Mentre ci pensavo su, comparve Diane sulla porta. «L'offerta è ancora
valida? Per il viaggio in Florida, dicevi sul serio?»

«Certo.»

«Perché ho parlato con Simon. Non è proprio felice all'idea ma per
qualche altro giorno starà bene anche da solo.»

``Veramente riguardoso da parte sua'' pensai.

«Quindi, a meno che\ldots{} lo so che avevi bevuto, ma\ldots»

«Non essere sciocca. Chiamo la compagnia aerea.»

Prenotai un posto a nome di Diane sul primo volo D.C/Orlando del giorno
dopo.

Poi terminai di fare i bagagli. Delle cose di mia madre, alla fine presi
i due Buddha di giada scheggiata.

Guardai in giro per la casa, controllai anche sotto i letti ma la
scatola mancante {RICORDI} ({SCUOLA}) sembrava svanita per sempre.